\section{Hilbert projective geometry and the exact distortion}
\label{sec:tpc49-Hilbert-geometry}

Let \(\Hh_1,\Hh_2\) be finite-dimensional Hilbert spaces and identify
\(T\in\Hh_1\otimes\overline{\Hh_2}\) with an operator
\(\Hh_2\to\Hh_1\).  Standard facts about projective tensor products,
nuclear operators, and Hilbert--Schmidt operators may be found in
\citep{Ryan2002,HornJohnson2013,Simon2005}; the short
proofs needed here are included.

\begin{definition}[Hilbert projective envelope]
\label{def:tpc49-Hilbert-envelope}
For a rank-one representation written as
\begin{equation}
 \mathscr R:\qquad
 T=\sum_{\omega\in\Omega}
 x_\omega\otimes\overline{y_\omega},
 \label{eq:tpc49-rank-one-representation}
\end{equation}
define its envelope and optimum by
\begin{equation}
 \cE_{\mathscr R}(T)
 :=\sum_{\omega\in\Omega}
 \|x_\omega\|_2\|y_\omega\|_2,
 \qquad
 \cE_{\rm pr}(T):=\inf_{\mathscr R}\cE_{\mathscr R}(T).
 \label{eq:tpc49-Hilbert-envelope}
\end{equation}
Finite sums suffice.  A Bochner integral gives the same infimum after
finite-dimensional approximation.
\end{definition}

\begin{theorem}[Hilbert projective equals nuclear]
\label{thm:tpc49-Hilbert-projective}
If \(T\ne0\) has singular values \(\sigma_1,\ldots,\sigma_r>0\), then
\begin{align}
 \cE_{\rm pr}(T)&=\|T\|_{S_1}=\sum_{k=1}^r\sigma_k,
 \label{eq:tpc49-projective-equals-nuclear}\\
 \cD_{\rm at}(T)&:=\|T\|_{S_2}^2=\sum_{k=1}^r\sigma_k^2,
 \label{eq:tpc49-atomic-is-Frobenius}\\
 \cE_{\rm pr}(T)^2
 &=r_{\rm nuc}(T)\cD_{\rm at}(T),
 \qquad
 r_{\rm nuc}(T)
 :=\frac{(\sum_k\sigma_k)^2}{\sum_k\sigma_k^2}.
 \label{eq:tpc49-nuclear-participation-rank}
\end{align}
Moreover
\begin{equation}
 1\le \operatorname{sr}(T)
 :=\frac{\|T\|_{S_2}^2}{\|T\|_{\rm op}^2}
 \le r_{\rm nuc}(T)
 \le\operatorname{rank}T.
 \label{eq:tpc49-effective-rank-chain}
\end{equation}
The upper equality holds exactly when all nonzero singular values are
equal; the lower endpoint \(r_{\rm nuc}=1\) holds exactly at rank one.
\end{theorem}

\begin{proof}
The singular-value decomposition
\(T=\sum_k\sigma_ku_k\otimes\overline{v_k}\) gives
\(\cE_{\rm pr}(T)\le\sum_k\sigma_k\).  Conversely, the nuclear norm is
a norm and every rank-one operator satisfies
\(\|x\otimes\overline y\|_{S_1}=\|x\|_2\|y\|_2\).  Thus every
representation \eqref{eq:tpc49-rank-one-representation} has
\(\|T\|_{S_1}\le\cE_{\mathscr R}(T)\), proving
\eqref{eq:tpc49-projective-equals-nuclear}.  Formula
\eqref{eq:tpc49-atomic-is-Frobenius} is the Hilbert--Schmidt identity.

Cauchy--Schwarz gives
\((\sum_k\sigma_k)^2\le r\sum_k\sigma_k^2\), with equality exactly
for equal positive singular values.  Also
\(\sigma_1\sum_k\sigma_k\ge\sum_k\sigma_k^2\), whence
\begin{equation}
 r_{\rm nuc}(T)
 \ge\frac{\sum_k\sigma_k^2}{\sigma_1^2}
 =\operatorname{sr}(T).
 \label{eq:tpc49-stable-lower-proof}
\end{equation}
The remaining equality statements are immediate.
\end{proof}

\begin{remark}[Stable rank points in the wrong direction]
\label{rem:tpc49-stable-rank-insufficient}
Stable rank is a lower, not an upper, bound for the present distortion.
Let \(T_N\) have singular values \(1\) and \(N^{-3/4}\), the latter
repeated \(N\) times.  Then
\begin{equation}
 \operatorname{sr}(T_N)=1+N^{-1/2}\longrightarrow1,
 \qquad
 r_{\rm nuc}(T_N)
 =\frac{(1+N^{1/4})^2}{1+N^{-1/2}}
 \asymp N^{1/2}.
 \label{eq:tpc49-stable-rank-counterexample}
\end{equation}
Thus the nuclear participation rank, not ordinary stable rank, is the
correct parameter.
\end{remark}

\subsection{A fixed representation has a second condition number}

\begin{definition}[Representation cancellation number]
\label{def:tpc49-representation-condition}
For a fixed exact representation \(\mathscr R\) of \(T\ne0\), define
\begin{equation}
 c_{\mathscr R}(T)
 :=\frac{\cE_{\mathscr R}(T)}{\|T\|_{S_1}}\ge1.
 \label{eq:tpc49-representation-condition}
\end{equation}
\end{definition}

\begin{proposition}[Exact fixed-representation ledger]
\label{prop:tpc49-fixed-representation-ledger}
For every fixed exact rank-one representation,
\begin{equation}
 \boxed{
 \frac{\cE_{\mathscr R}(T)^2}{\cD_{\rm at}(T)}
 =c_{\mathscr R}(T)^2r_{\rm nuc}(T).}
 \label{eq:tpc49-fixed-representation-ledger}
\end{equation}
There is no finite upper bound for the left side that is uniform over
all exact signed representations of one fixed nonzero tensor.
\end{proposition}

\begin{proof}
The identity is the product of the definitions.  For nonuniformity,
fix a nonzero rank-one tensor \(z=x\otimes\overline y\), normalized so
that \(\|x\|\|y\|=1\).  For every \(R>0\), the identity
\begin{equation}
 z=(R+1)z-Rz
 \label{eq:tpc49-cancelling-pair}
\end{equation}
has envelope \(2R+1\), although the represented tensor, its rank, and
its atomic energy remain unchanged.  More generally, the null pair
\(Rz-Rz\) can be appended to any representation.
\end{proof}

This elementary obstruction is consequential for
\eqref{eq:tpc49-inherited-envelope}.  A theorem for a specified
TPC representation may estimate its cancellation number.  A pointwise
infimum may suppress artificial null pairs.  A claim for every exact
signed representation cannot be true.

\subsection{Two usable positive criteria}

\begin{corollary}[Approximate nuclear rank]
\label{cor:tpc49-approximate-rank}
Suppose the singular values are decreasing and, for some \(R\ge1\),
\begin{equation}
 \sum_{k>R}\sigma_k\le\eta\|T\|_{S_2}.
 \label{eq:tpc49-nuclear-tail}
\end{equation}
Then
\begin{equation}
 \|T\|_{S_1}
 \le(\sqrt R+\eta)\|T\|_{S_2},
 \qquad
 r_{\rm nuc}(T)\le(\sqrt R+\eta)^2.
 \label{eq:tpc49-approximate-rank-bound}
\end{equation}
\end{corollary}

\begin{proof}
Cauchy--Schwarz bounds the first \(R\) singular values by
\(\sqrt R\|T\|_{S_2}\); add \eqref{eq:tpc49-nuclear-tail}.
\end{proof}

\begin{proposition}[Orthogonal layers do not improve an \(L^1\) envelope]
\label{prop:tpc49-orthogonal-layers}
Suppose the rank-one tensors \(z_1,\ldots,z_s\) are pairwise orthogonal
in Hilbert--Schmidt inner product, put \(a_k=\|z_k\|_{S_2}\), and use
the displayed decomposition \(T=\sum_kz_k\).  Then
\begin{equation}
 \cD_{\rm at}(T)=\sum_ka_k^2,
 \qquad
 \cE_{\mathscr R}(T)^2
 =N_{\rm eff}(a)\cD_{\rm at}(T),
 \qquad
 N_{\rm eff}(a):=\frac{(\sum_ka_k)^2}{\sum_ka_k^2}.
 \label{eq:tpc49-orthogonal-layer-loss}
\end{equation}
Equal layers lose the full layer count \(s\).
\end{proposition}

Thus orthogonality helps only if it is retained through an
\(L^2\)/Bessel reassembly.  Applying Minkowski first converts it into
an \(L^1\) layer count.  Entrywise positivity or one-sign coefficients
do not repair this: permutation matrices below are nonnegative and
still attain maximal nuclear distortion.
